\section{関連研究}
\subsection{閉域的な物体検出モデル}
YOLO\cite{ref_YOLO}やDETR\cite{ref_DETR}に代表される物体検出モデルの多くは，ImageNet-1k\cite{ref_ImageNet}やCOCO\cite{ref_COCO}のように，あらかじめ定めた固定のクラス集合を持つデータセットを用いて学習される．
特にDETRは，Transformerに基づくEnd-to-Endの物体検出モデルであり，候補領域の生成やNon-Maximum Suppression（NMS）といった後処理を必要とせず，集合予測に基づく損失関数により学習される．
これらのモデルは，学習時に定めたクラス集合の範囲内では高い精度を示す一方，学習時に存在しなかったクラスの物体を原理的に予測できない．
本研究では，YOLOの近年の実装であるYOLO11およびYOLOv8\cite{ref_Ultralytics}を比較対象とし，4.8節で提案手法との精度の差と，その差が検出器と分類器のいずれに起因するかを評価する．
本研究では，DETRの予測カテゴリそのものではなく，そのバウンディングボックス出力のみを物体領域の検出に用いる．
すなわち，DETRが学習していない未知のカテゴリの物体であっても，その物体領域さえ検出できれば，後段のVLMによって開語彙のカテゴリをアノテーションできる．

\subsection{オープンボキャブラリ物体検出モデル}
OWL-ViT\cite{ref_OWLViT}やGrounding DINO\cite{ref_GDINO}は，候補となるカテゴリ名や説明文をテキストクエリとして与えることで，学習時に固定されたクラス集合に依存せず物体を検出する．
これらのモデルは物体領域の検出とカテゴリの予測を単一のモデルで同時に行うため，候補カテゴリ集合を変更するたびに検出処理を再実行する必要がある．
候補カテゴリ集合がWordNetの階層のように多段階かつ大規模である場合，例えば\texttt{person.n.01}の直接の下位語だけで約400語あり，各段階・各分岐でこの再検出を繰り返すことは計算コストの観点から非効率である．
本研究の設計は，物体領域の検出を1回だけ実行し，検出済みの領域に対する軽量なVLMの推論のみを繰り返す点で，これらのモデルと異なる．

\subsection{視覚言語モデルによるゼロショット分類}
CLIP\cite{ref_CLIP}は，画像とその説明文のペアから構成される大規模データセットを用い，対照学習により画像エンコーダとテキストエンコーダを同時に学習するモデルである．
学習後のCLIPは，任意の単語や文をテキストエンコーダに入力し，画像エンコーダの出力との類似度を計算することで，学習時に固定のクラス集合を定めることなく画像分類（ゼロショット分類）を行うことができる．
本研究では，このゼロショット分類の性質を利用し，検出済みの物体領域のクロップに対して直接カテゴリのアノテーションを行う．

\subsection{階層的・粒度可変な分類}
ImageNetの多くのクラスはWordNet\cite{ref_WordNet}の同義語集合（Synset）に対応付けられており，画像分類においてクラス階層を利用する研究は古くから存在する．
しかし，従来の階層的画像分類の多くは，あらかじめ定めた固定の階層構造に対して学習を行うものであり，推論時に階層の深さや分岐を動的に選択することは想定していない．
本研究の粒度可変アノテーションは，学習を一切行わず，推論時にVLMへ与える候補語彙集合を動的に変更するだけで，任意の深さまでの粒度でアノテーションを行える点で異なる．
